% !TEX encoding = UTF-8 Unicode
% !TEX TS-program = XeLaTeX
% !TEX root = ParliamentUserGuide.tex

%%%%%%%%%%%%%%%%%%%%%%%%%%%%%%%%%%%%%%%%%%%%%%%%%%
\chapter{Deploying and Using \acl{pmnt}}
\label{chapter-deploying-and-using}

\ac{pmnt} is a server application that exposes a \ac{sparql} endpoint.  \ac{pmnt} can be deployed either using its Docker image (see Section~\ref{section-std-container-deploy}) or a more traditional binary distribution (see Section~\ref{section-std-server-deploy}). Pre-built binaries for common platforms are available on the \ac{pmnt} Web site\urlcite{https://github.com/\githubOrg/parliament/releases}.  You can also build \ac{pmnt} yourself (see Chapter~\ref{chapter-building-parliament}).

In addition to the deployment instructions in Sections~\ref{section-std-container-deploy} and~\ref{section-std-server-deploy}, this chapter discusses \ac{pmnt} configuration files (Section~\ref{section-configuring-pmnt}), \ac{pmnt} utilities (Section~\ref{section-parliamentadmin-utility}), and troubleshooting (Section~\ref{section-troubleshooting}).

%%%%%%%%%%%%%%%%%%%%%%%%%%%%%%%%%%%%%%%%%%%%%%%%%%
\section{Deploying a \ac{pmnt} Container}
\label{section-std-container-deploy}

The following command will run the container (line breaks added for readability). Table~\ref{tbl:DockerTags} lists the Docker Hub tags of \acp{pmnt} Docker images.
{\small
\begin{verbatim}
docker run -d -p 80:8089 --name=parliament-«version»
   --mount type=volume,src=kb-data,dst=/var/parliament-data
   --env PARLIAMENT_JAVA_HEAP_SIZE=1g
   idemmons/parliament-«arch»:«version»
\end{verbatim}
}

\begin{table}[htbp]
	\centering
	\begin{tabular}{ll}
		\toprule
		Architecture & Tag\\
		\headingrule
		Intel and AMD 64-bit: & \texttt{idemmons/parliament-x64:«version»}\\
		ARM-64: & \texttt{idemmons/parliament-a64:«version»}\\
		\bottomrule
	\end{tabular}
	\caption{\acp*{pmnt} Docker tags}
	\label{tbl:DockerTags}
\end{table}

As given here, the mount option will create a local Docker volume called \path|kb-data| to contain \acp{pmnt} data files and logs.  Using variations of this option, you can cause the container to store the data files and logs in a broad variety of locations and file systems, such as on an external storage array.  (See the Docker documentation for more information.)

Another common storage option places \acp{pmnt} data files in a local directory on the container's host like so:
{\small
\begin{verbatim}
docker run -d -p 80:8089 --name=parliament-«version»
   --mount type=bind,src=/host/directory,dst=/var/parliament-data
   --env PARLIAMENT_JAVA_HEAP_SIZE=1g
   idemmons/parliament-«arch»:«version»
\end{verbatim}
}

The environment variable in this command is optional. If it is left out, the default value used will be the value shown in the command lines above (see Section~\ref{section-pmnt-startup-script}). Another environment variable that may be useful here is \verb|PARLIAMENT_KB_CONFIG_PATH| (see section~\ref{section-main-config}).

Note that the same effects as the above two command lines can be achieved through the Docker GUI as well. First download the image using its tag, and then press the Run button. Click Optional Settings, and then fill in the container name, port, volume information, and environment variable(s).

See Table~\ref{table-pmnt-connect-urls} in Section~\ref{section-std-server-usage} for a list of the endpoints exposed by the running container and a discussion of how to use them.



%%%%%%%%%%%%%%%%%%%%%%%%%%%%%%%%%%%%%%%%%%%%%%%%%%
\section{Deploying a \ac{pmnt} Server}
\label{section-std-server-deploy}

A \ac{pmnt} distribution is a compressed archive containing a ready-to-run \ac{pmnt} server, with the following naming convention:
\begin{quote}
	\texttt{parliament-\textbf{\textit{x.y.z}}-\textbf{\textit{platform}}-\textbf{\textit{architecture}}.zip}
\end{quote}
Here \texttt{\textbf{\textit{x.y.z}}} is the \ac{pmnt} version number, \texttt{\textbf{\textit{platform}}} indicates the platform (see Table~\ref{tbl:SupportedPlatforms}), and \texttt{\textbf{\textit{architecture}}} indicates the hardware architecture (see Table~\ref{tbl:SupportedArchitectures}).

\begin{table}[htbp]
	\centering
	\begin{tabular}{ll}
		\toprule
		\emph{Platform}	& \emph{Description}				\\
		\headingrule
		\path|mac|			& Mac OS Catalina and newer	\\
		\path|ubuntu26|	& Ubuntu Linux 26, LTS			\\
		\path|win|			& Windows 11						\\
		\bottomrule
	\end{tabular}
	\caption{Supported Platforms}
	\label{tbl:SupportedPlatforms}
\end{table}

\begin{table}[htbp]
	\centering
	\begin{tabular}{ll}
		\toprule
		\emph{Architecture}	& \emph{Description}				\\
		\headingrule
		\path|-x64|				& Intel and AMD 64-bit			\\
		\path|-a64|				& ARM-64								\\
		\emph{Blank}			& Universal binary (Mac only)	\\
		\bottomrule
	\end{tabular}
	\caption{Supported Hardware Architectures}
	\label{tbl:SupportedArchitectures}
\end{table}

If you are building \ac{pmnt} yourself according to the instructions in Chapter~\ref{chapter-building-parliament}, then these archives appear in the \path|target| directory at the end of the build process.

\subsection{Installing for the First Time}
\label{section-std-server-init-deploy}

Once you have downloaded a distribution, these steps will yield a functioning \ac{kb}:
\begin{enumerate}
	\item Extract the archive to your preferred location.  On UNIX-like systems, \path|/usr/opt/ParliamentKB| or \path|/usr/local/ParliamentKB| is a likely choice.  On Windows, \path|C:\Program Files\ParliamentKB| is customary.  The instructions that follow refer to this directory as simply \path|ParliamentKB|.

	\item On Windows, install the C and C++ run-time libraries.  You can find installers for these in the following directory:
	\begin{quote}
		\path|ParliamentKB/RedistributablePackages|
	\end{quote}

	\item On Windows, open a PowerShell, switch to the \path|ParliamentKB| directory, and invoke the startup script to start the \ac{pmnt} server:
	\begin{quote}
		\path|parliament.ps1|
	\end{quote}
	If you prefer the Command Prompt, type:
	\begin{quote}
		\path|parliament.bat|
	\end{quote}

	On other platforms, open a shell, switch to the \path|ParliamentKB| directory, and start \ac{pmnt} using this command:
	\begin{quote}
		\texttt{./parliament}
	\end{quote}

	Whatever platform you are using, this starts \ac{pmnt} as a command line process running under the current terminal.

	\item Shut down the server by typing ``exit'' or ``Ctrl+C''.
\end{enumerate}



\subsection{Upgrading an Existing Installation}
\label{section-std-server-upgrade}

When upgrading an existing \ac{pmnt} installation to a newer version, the procedures of Section~\ref{section-std-server-init-deploy} are not sufficient because they ignore the migration of existing data.  The \ac{pmnt} development team strives to maintain compatibility of file formats whenever possible.  In such cases, all that is required is to move the \path|kb-data| subdirectory of your \ac{pmnt} installation into the new \ac{pmnt} installation.  If you have customized \acp{pmnt} configuration, you will also need to make those changes in the new installation as well.

However, from time to time \acp{pmnt} file formats do change, and in such cases a different migration procedure is required.  Following this procedure even in those cases where it is not required has some benefits as well, since it frees up any unused space in the \ac{pmnt} data files.  This is the way to perform such a migration:

\begin{enumerate}
	\item\label{upgrade-procedure-first-step}\textbf{\emph{Before upgrading,}} point your browser at your \ac{pmnt} server.  The \ac{url} for this is \nolinkurl{http://<host>:<port>/parliament/}.  By default this will be \nolinkurl{http://localhost:8089/parliament/}.

	\item Click on the ``Export Data'' link.

	\item\label{item-export-download}Click on the ``Backup Repository'' button.  This will start a download whose file name looks like this:
\begin{quote}\path|parliament-export-localhost-<date>-<time>.zip|\end{quote}

	\item When the download finishes, shut down \ac{pmnt}.  Rename the \ac{pmnt} installation directory if you want your upgraded installation to reside in the same place.

	\item Unzip the new version of \ac{pmnt} to your desired location, customize the new installation's settings as required (see Section~\ref{section-configuring-pmnt}), and start up the new \ac{pmnt} installation.

	\item Point your browser at the new \ac{pmnt} server.  The \ac{url} is the same as the one in Step~\ref{upgrade-procedure-first-step} above.

	\item Click on the ``Insert Data'' link.  In the ``Restore Repository Backup'' section, press the ``Choose File'' button and select the file you downloaded in Step~\ref{item-export-download} above.  Then press the ``Restore Repository Backup'' button.
\end{enumerate}

Once you see the ``Import Complete'' message, your upgraded \ac{pmnt} server is ready for business.

\subsection{Usage}
\label{section-std-server-usage}

\subsubsection{Shutting \ac{pmnt} Down}

\emph{Important note:}  When you shut down \ac{pmnt}, it is important to shut it down gracefully.  Otherwise, there is a small probability that the \ac{pmnt} files may not be flushed to disk before they are closed, and they may be corrupted.  If you run \ac{pmnt} as a Container (see Section~\ref{section-std-server-init-deploy}), this is generally not a problem, because the operating system sends a shutdown message at the appropriate times.  However, if you run \ac{pmnt} from the command line, you will need to shut it down yourself by typing ``exit'' or ``Ctrl+C''.

\subsubsection{\ac{pmnt} Connection \acp{url}}

To use the \ac{pmnt} server, you will need the connection \acp{url} listed in Table~\ref{table-pmnt-connect-urls}. Note that these need to be prefixed with \nolinkurl{http://«host»:«port»} (or \nolinkurl{http://localhost:8089} by default).  See Section~\ref{section-configuring-pmnt} to configure your server to run on a host name and port other than ``localhost'' and ``8089.''
\begin{table}[htbp]
	\centering\small
	\begin{tabular}{ll}
		\toprule
		Endpoint Description & \ac{url} Path\\
		\headingrule
		Web interface: & \nolinkurl{/parliament}\\
		\ac{sparql} endpoint: & \nolinkurl{/parliament/sparql}\\
		\ac{sparql} Update endpoint: & \nolinkurl{/parliament/update}\\
		\ac{sparql} Graph Store \acs*{http} Protocol: & \nolinkurl{/parliament/graphstore}\\
		Legacy bulk insert/export endpoint: & \nolinkurl{/parliament/bulk}\\
		\bottomrule
	\end{tabular}
	\caption{\acs*{pmnt} Server Connection \acsp*{url}}
	\label{table-pmnt-connect-urls}
\end{table}

\subsubsection{\acp{pmnt} Web Interface}

You can use the \ac{pmnt} server interactively by pointing a web browser at the first \ac{url} in Table~\ref{table-pmnt-connect-urls}.  This simple and (hopefully) self-explanatory web interface is useful for initial data loading, experimenting with queries, exploring a collection of data, troubleshooting, and similar tasks.

\subsubsection{Using \ac{pmnt} Programmatically}

In order to issue queries programmatically, you will need to write some code to send \ac{sparql}-compliant \ac{http} requests.  One library that can send such requests on your behalf is Jena itself.  Figure~\ref{figure-sparql-select-query-with-jena} shows sample code for issuing a \ac{sparql} select query.
\begin{figure}[htbp]
	\footnotesize
	\VerbatimInput{includes/jena-select.java}
	\caption{Issuing a \acs*{sparql} select query with Jena}
	\label{figure-sparql-select-query-with-jena}
\end{figure}
Note that this code uses the \ac{sparql} endpoint \ac{url} from Table~\ref{table-pmnt-connect-urls}.

Similarly, we can also update the content of \ac{pmnt} programmatically using the Jena library to send \ac{sparql}-Update requests.  Figure~\ref{figure-sparql-update-with-jena} shows sample code for issuing an update.
\begin{figure}[htbp]
	\footnotesize
	\VerbatimInput{includes/jena-update.java}
	\caption{Issuing a \acs*{sparql}-Update request with Jena}
	\label{figure-sparql-update-with-jena}
\end{figure}

Another common update operation is inserting a body of \ac{rdf} into \ac{pmnt}, shown in Figure~\ref{figure-sparql-insert-with-jena}.
\begin{figure}[htbp]
	\footnotesize
	\VerbatimInput{includes/jena-insert.java}
	\caption{Inserting \acs*{rdf} with Jena}
	\label{figure-sparql-insert-with-jena}
\end{figure}

\subsubsection{Inserting Data into \ac{pmnt} from the Command Line}

The \path|curl| command line utility can insert an \ac{rdf} file into a running \ac{pmnt} instance via the \ac{sparql} Graph Store \acs*{http} Protocol (line breaks inserted for readability):

\begin{verbatim}
  curl -H 'Content-Type: text/turtle' --data-binary '@file.ttl'
    'http://localhost:8089/parliament/graphstore?default'
\end{verbatim}

If the command succeeds, \path|curl| will print a response containing the strings ``HTTP OK: 200'' and ``Inserted \emph{N} statements.''

To insert into a named graph, replace the \path|?default| portion of the \ac{url} with \path|?graph=«graph-iri»|, where \path|«graph-iri»| is the \ac{iri} of the named graph. For example:

\begin{verbatim}
  curl -H 'Content-Type: text/turtle' --data-binary '@file.ttl'
    'http://localhost:8089/parliament/graphstore
    ?graph=http://example.org/graph1'
\end{verbatim}

This will create the named graph if it does not already exist.



\subsection{Securing a \ac{pmnt} Server}
\label{section-securing-parliament}

This section discusses how to use an Apache web server (httpd) or NGINX on Linux as a secure proxy for a \ac{pmnt} server running on the same machine.  This procedure will implement TLS so that the communication channel is encrypted and the server is authenticated to clients.  Client authentication is not covered at this time.

Before we begin, update the operating system.  (Note that the commands below use \texttt{yum}.  Please substitute \texttt{apt} for \texttt{yum} if your Linux distribution uses that package manager instead.)

\begin{Verbatim}
	sudo yum update
\end{Verbatim}

Next install required software.

\begin{Verbatim}
	sudo yum install «java-development-kit»
\end{Verbatim}

where \texttt{«java-development-kit»} is \texttt{java-8-amazon-corretto-headless} on AWS and \texttt{openjdk-8-jdk-headless} in other environments.  Next install your desired Web server to serve as a proxy.  To use Apache:

\begin{Verbatim}
	sudo yum install httpd mod_ssl
\end{Verbatim}

Or to use NGINX:

\begin{Verbatim}
	sudo yum install nginx-light
\end{Verbatim}

On your virtual machine, configure the following:

\begin{itemize}
	\item Deploy \ac{pmnt} in the standard way.  (See Section~\ref{section-std-server-init-deploy} or~\ref{section-std-server-upgrade} above for details.)

	\item It is crucial that in the startup script \texttt{parliament}, located in the root directory of your \ac{pmnt} installation, the JETTY\_HOST setting near the top \emph{must} be set to localhost (or 127.0.0.1).  This prevents \ac{pmnt} itself from being directly exposed to network traffic from outside the machine.

	\item In those same two files, ensure JETTY\_PORT is set to 8089.  In reality, most any port may be used here.  8089 is \acp{pmnt} default, and a useful choice.  The instructions below assume this port.

	\item Ensure that any firewalls (either on the machine itself or elsewhere) are set to allow HTTP (port 80) and HTTPS (port 443) traffic to reach the machine from any network locations that need access to the \ac{pmnt} server.
\end{itemize}

Acquire a web server certificate for the machine from your favorite certificate authority.  This typically includes three files:

\begin{itemize}
	\item The certificate itself, referred to as \texttt{certificate.crt} below.

	\item The private key, called \texttt{private.key} below.

	\item The root and intermediate certificate chain, called \texttt{ca-bundle.crt} below.
\end{itemize}



\subsubsection{Securing \ac{pmnt} with Apache}
\label{section-securing-parliament-apache}

Create a directory for the three certificate files and store them there.  The actual location is not important, but we will refer to it as \texttt{/etc/server-certs} below.  The three files, and especially the private key, should be stored only at this location on the server to prevent unauthorized disclosure.  A backup on a separate, secure machine is also a good idea.  To finish securing the files, run these commands:

\begin{Verbatim}
	sudo chmod u=rw,go= /etc/server-certs/*
	sudo chown root.root /etc/server-certs/*
\end{Verbatim}

Copy the following file from your \ac{pmnt} installation's \texttt{conf} directory to \texttt{/etc/httpd/conf.d}:

\begin{Verbatim}
	parliament-redirect-httpd.conf
\end{Verbatim}

In that file, substitute your machine's host name (the host name that you used when you requested the certificate) in the indicated places, and be sure that the file contains the correct paths to the three certificate files and the correct port for \ac{pmnt}.  Then restart the web server to activate the new configuration with this command:

\begin{Verbatim}
	sudo systemctl restart httpd
\end{Verbatim}



\subsubsection{Securing \ac{pmnt} with NGINX}
\label{section-securing-parliament-nginx}

Create a directory for the three certificate files and store them there.  The actual location is not important, but we will refer to it as \texttt{/etc/server-certs} below.  The three files, and especially the private key, should be stored only at this location on the server to prevent unauthorized disclosure.  A backup on a separate, secure machine is also a good idea.

Unlike Apache, NGINX requires that the certificate and certificate chain be combined into a single file, like so:

\begin{Verbatim}
	cat certificate.crt ca-bundle.crt > temp.crt
	rm certificate.crt ca-bundle.crt
	mv temp.crt certificate.crt
\end{Verbatim}

Note that the order in which the two files are combined is important --- the certificate must come before the certificate chain.  Finish securing the files with these commands:

\begin{Verbatim}
	sudo chmod u=rw,go= /etc/server-certs/*
	sudo chown root.root /etc/server-certs/*
\end{Verbatim}

Copy the following file from your \ac{pmnt} installation's \texttt{conf} directory to \texttt{/etc/nginx/sites-available}:

\begin{Verbatim}
	parliament-redirect-nginx.conf
\end{Verbatim}

In that file, substitute your machine's host name (the host name that you used when you requested the certificate) in the indicated places, and be sure that the file contains the correct paths to the two certificate files and the correct port for \ac{pmnt}.  From \texttt{/etc/nginx/sites-enabled}, create a symbolic link:

\begin{Verbatim}
	ln -s /etc/nginx/sites-available/parliament-redirect-nginx.conf
\end{Verbatim}

Finally, restart the web server to activate the new configuration with this command:

\begin{Verbatim}
	sudo systemctl restart nginx
\end{Verbatim}



%%%%%%%%%%%%%%%%%%%%%%%%%%%%%%%%%%%%%%%%%%%%%%%%%%
\section{Configuring \ac{pmnt}}
\label{section-configuring-pmnt}

The section begins with a how-to for common configuration tasks, followed by a more detailed discussion of each of \acp{pmnt} configuration files.

%%%%%%%%%%%%%%%%%%%%%%%%%%%%%%%%%%%%%%%%%%%%%%%%%%
\subsection{Common Configuration Tasks}
\label{section-common-config-tasks}

\emph{Changing the host name and port:} See Section~\ref{section-pmnt-startup-script}.

\emph{Changing the location of \acp{pmnt} data files:} See the \texttt{kbDirectoryPath} setting in Section~\ref{section-main-config}.

\emph{Setting the query timeout:} See the \texttt{TimeoutDuration} and \texttt{TimeoutUnit} settings in Section~\ref{section-main-config}.

\emph{Memory configuration:} \ac{pmnt} uses three separate pools of memory:
\begin{itemize}[noitemsep]
	\item The Java heap (to configure, see Section~\ref{section-pmnt-startup-script})
	\item The demand-paged cache maintained by the operating system for memory-mapped files (not configurable)
	\item The C++ heap (not configurable)
\end{itemize}

Achieving the right balance between these pools of memory is tricky, but the key thing to keep in mind is, don't make the Java heap too large. Java programmers are accustomed to setting their Java heap as large as possible, but for \ac{pmnt}, the Java heap is not as important as the memory-mapped files. The Java heap size should be significantly less than the total memory available to leave sufficient memory for the memory-mapped files.

\emph{Concurrency:} \ac{pmnt} has a many readers/single writer model of concurrency.  In other words, if \ac{pmnt} receives several requests at once, and all of them are read-only requests (such as a query), then \ac{pmnt} will execute those requests in parallel.  Whenever \ac{pmnt} receives a request that requires writing (such as an insert, a \ac{sparql} Update, or a graph creation), that request will wait until all other executing requests are finished, and then it will execute. And while it is executing, all other requests (read or write) will wait in the queue until it is finished.

This model works well in most situations, where the number of concurrent requests is not terribly large and read-only requests dominate.  A few dozen concurrent queries is not uncommon in our experience, and works well.  Larger numbers of concurrent read-only requests will work correctly, but the performance may degrade because all those requests ultimately are reading from the same I/O device.  Therefore, caching becomes the key issue for concurrency levels this high (see memory configuration above).

%%%%%%%%%%%%%%%%%%%%%%%%%%%%%%%%%%%%%%%%%%%%%%%%%%
\subsection{\acp{pmnt} Startup Script}
\label{section-pmnt-startup-script}

\acp{pmnt} startup script, located in the root installation directory, contains a few configuration settings.  On Windows it is the PowerShell script \path|parliament.ps1|, while other platforms use the bash script \path|parliament|.

Change the host name (or IP address) and port of the \ac{sparql} endpoint by setting \verb|SPRING_BOOT_HOST| and \verb|SPRING_BOOT_PORT| in the startup script (\verb|$springBootHost| and \verb|$springBootPort| on Windows).  These default to \verb|localhost| and 8089, respectively.  To make \ac{pmnt} addressable on any of a host's network interfaces, change \verb|SPRING_BOOT_HOST| to ``0.0.0.0''.

Setting \verb|DEFAULT_JAVA_HEAP_SIZE| (\verb|$defaultJavaHeapSize| on Windows) controls the amount of memory set aside for the Java heap. Section~\ref{section-common-config-tasks} above offers guidelines for this setting. The syntax for this setting is an integer concatenated with an `m' (MB) or `g' (GB), and the default is `1g'.

%%%%%%%%%%%%%%%%%%%%%%%%%%%%%%%%%%%%%%%%%%%%%%%%%%
\subsection{\acp{pmnt} Main Configuration File}
\label{section-main-config}

\acp{pmnt} main configuration file contains settings pertaining to storage, query, and inference.  It is typically called \path|ParliamentKbConfig.txt|, and \ac{pmnt} finds it like so:

\begin{enumerate}
	\item If the environment variable \verb|PARLIAMENT_KB_CONFIG_PATH| is set, its value is assumed to be the configuration file's absolute path.  Note that using this option, you can rename the file to anything you wish.

	\item If this environment variable is not set, then \ac{pmnt} looks for a file named \path|ParliamentKbConfig.txt| in the same directory as the \ac{pmnt} \ac{dll} (on Windows) or shared library (on other platforms).

	\item If that file is not found, then \ac{pmnt} looks for a file of the same name in the parent directory of the \ac{pmnt} \ac{dll} or shared library.  In practice, this is typically how the configuration file is found.

	\item Otherwise, \ac{pmnt} looks for a file named \path|ParliamentKbConfig.txt| in the current working directory of the server process.
\end{enumerate}

The configuration file is a plain text file containing name-value pairs.  The names are case-insensitive.  Blank lines and comment lines (preceded by a `\#') are ignored.  Many of the settings are Boolean values, in which case the value is case-insensitive and may be ``true'', ``t'', ``yes'', ``y'', ``on'', or ``1'' (for true), or ``false'', ``f'', ``no'', ``n'', ``off'', or ``0'' (for false).

The recognized settings are as follows:
\begin{description}
	\item[kbDirectoryPath] The path of the \ac{pmnt} knowledge base files.  If they do not yet exist, they will be created here.  Log files will also be stored here in the \path|log| sub-directory.  \emph{Default: \path|kb-data|}

	\item[stmtFileName] The name of the memory-mapped file containing records of statements. \emph{Default: \path|statements.mem|}

	\item[rsrcFileName] The name of the memory-mapped file that contains resource records. \emph{Default: \path|resources.mem|}

	\item[uriTableFileName] The name of the memory-mapped file containing resource strings (URIs and literals). \emph{Default: \path|uris.mem|}

	\item[uriToIntFileName] The name of the directory containing the mapping from numeric resource identifiers to resource strings.  The contents of this directory are managed by RocksDB. \emph{Default: \path|u2i.db|}

	\item[stmtToIdFileName] This obsolete setting is now ignored.

	\item[keepDupStmtIdx] This obsolete setting is now ignored.

	\item[readOnly] When set to ``yes'', prevents \ac{pmnt} from changing the underlying storage files in any way. \emph{Default: ``no''}

	\item[fileSyncTimerDelay] \ac{pmnt} periodically flushes its underlying data files to disk to decrease the chances of a file corruption and to limit the amount of time required to shut down \ac{pmnt} gracefully.  This parameter is the time interval in milliseconds between flushes.  Set it to zero to disable flushing the files to disk.  This setting applies only to \ac{pmnt} deployed as a server.  Any other deployment will ignore this setting. \emph{Default: ``15000''}

	\item[initialRsrcCapacity] The number of resources \ac{pmnt} should allocate space for when creating a new \ac{kb}. \emph{Default: ``300000''}

	\item[avgRsrcLen] The average resource length (in characters) that \ac{pmnt} should anticipate when allocating space in a new \ac{kb}. \emph{Default: ``100''}

	\item[rsrcGrowthIncrement] The number of resources by which \ac{pmnt} increases the resource table size when it runs out of space in the file.  If this is less than one, then linear growth of the resource table is disabled, in which case rsrcGrowthFactor must be larger than one (see below). \emph{Default: ``600000''}

	\item[rsrcGrowthFactor] The factor by which \ac{pmnt} increases the size of the resource table when it runs out of space in the file.  If this is less than or equal to one, then geometric growth of the resource table is disabled, in which case rsrcGrowthIncrement must be at least one (see above).  Note that this is a decimal number that should be formatted according to your locale, e.g., ``1.1`` in the US or ``1,1'' in Europe. \emph{Default: ``0''}

	\item[initialStmtCapacity] The number of statements \ac{pmnt} should allocate space for when creating a new \ac{kb}. \emph{Default: ``500000''}

	\item[stmtGrowthIncrement] The number of statements \ac{pmnt} adds to the statement table size when it runs out of space in the file.  If this is less than one, then linear growth of the statement table is disabled, in which case stmtGrowthFactor must be larger than one (see below). \emph{Default: ``1000000''}

	\item[stmtGrowthFactor] The factor by which \ac{pmnt} increases the size of the statement table when it runs out of space in the file.  If this is less than or equal to one, then geometric growth of the statement table is disabled, in which case stmtGrowthIncrement must be at least one (see above).  Note that this is a decimal number that should be formatted according to your locale, e.g., ``1.1`` in the US or ``1,1'' in Europe. \emph{Default: ``0''}

	\item[bdbCacheSize] This obsolete setting is now ignored.

	\item[normalizeTypedStringLiterals] Instructs \ac{pmnt} to convert typed string literals to plain literals, both upon insert and at query time, as mandated in \ac{rdf} 1.1.  In \ac{pmnt} versions before the implementation of this option (versions~2.7.9 and prior), the behavior was as if the option were set to ``no''.  Therefore, if you are migrating a \ac{pmnt} installation originally created with one of these older versions to a newer version, you must do one of two things:
	\begin{enumerate}
		\item Either before or after the upgrade, create an export of the entire triple store.  To do this, use the ``Export Repository'' option on the Export page of \acp{pmnt} Web-based administrative interface.  Then shut down \ac{pmnt} and delete its \path|kb-data| (or \path|data|) directory.  This is the location identified by the \texttt{kbDirectoryPath} setting in the configuration file.  Then, after upgrading, import the backup using the ``Import Repository'' option on the Insert Data page of the administrative interface.

		\item Set this option to ``no''.
	\end{enumerate}
	The first option is highly recommended.  It requires some effort, but will bring your triple store into compliance with \ac{rdf} 1.1 and deliver slightly better performance to boot.  \emph{Default: ``yes''}

	\item[TimeoutDuration] Sets the query execution timeout. \emph{Default: ``5''}

	\item[TimeoutUnit] Sets the units of the query execution timeout.  Valid values are ``nanoseconds'', ``microseconds'', ``milliseconds'', ``seconds'', ``minutes'', ``hours'', and ``days''. \emph{Default: ``minutes''}

	\item[runAllRulesAtStartup] Causes \ac{pmnt} to evaluate each of the enabled rules at startup against the existing statements in the \ac{kb} to see if any new statements should be materialized.  Generally, this is unnecessary, because the rules are always evaluated against each statement as it is asserted.  However, if some of the rules were initially disabled when statements were being asserted, and are then enabled, it may be necessary to turn this on for one startup of \ac{pmnt} to ensure that the state of the \ac{kb} is in sync with the rule settings.  Alternately, in such a case this setting may be left off and the ParliamentAdmin tool may be run against the \ac{kb} with the ``guaranteeEntailments'' option.  Note that when this setting is turned on, \ac{kb} the startup time may be lengthy. \emph{Default: ``no''}

	\item[enableSWRLRuleEngine] Enables an additional rule engine that implements the \ac{swrl} rule language. \emph{Default: ``no''}

	\item[SubclassRule] Enables the following inference rules:
\[A \subset B \land B \subset C \Rightarrow A \subset C\]
\[X \in A \land A \subset B \Rightarrow X \in B\]
\[X \in \text{rdfs:Class} \Rightarrow X \subset X\]
\[X \in \text{owl:Class} \Rightarrow X \subset X\]
where $\subset$ means ``sub-class of'' and $\in$ means ``of type''. \emph{Default: ``on''}

	\item[inferRdfsClass] When both this setting and ``SubclassRule'' are turned on, this enables the following inference rule:
\[A \subset B \Rightarrow A \in \text{rdfs:Class} \land B \in \text{rdfs:Class}\]
where $\subset$ means ``sub-class of'' and $\in$ means ``of type''. \emph{Default: ``off''}

	\item[inferOwlClass] When both this setting and ``SubclassRule'' are turned on, this enables the following inference rule:
\[A \subset B \Rightarrow A \in \text{owl:Class} \land B \in \text{owl:Class}\]
where $\subset$ means ``sub-class of'' and $\in$ means ``of type''. \emph{Default: ``off''}

	\item[inferRdfsResource] When both this and the ``SubclassRule'' settings are turned on, the following inference rule is enabled:
\[A \subset B \Rightarrow A \subset \text{rdfs:Resource} \land B \subset \text{rdfs:Resource}\]
where $\subset$ means ``sub-class of''. \emph{Default: ``off''}

	\item[inferOwlThing] When both this setting and ``SubclassRule'' are turned on, this enables the following inference rule:
\[A \subset B \Rightarrow A \subset \text{owl:Thing} \land B \subset \text{owl:Thing}\]
where $\subset$ means ``sub-class of''. \emph{Default: ``off''}

	\item[SubpropertyRule] Enables the following inference rules:
\[P \sqsubset Q \land Q \sqsubset R \Rightarrow P \sqsubset R\]
\[P \sqsubset Q \land P(X, Y) \Rightarrow Q(X, Y)\]
where $\sqsubset$ means ``sub-property of''. \emph{Default: ``on''}

	\item[DomainRule] Enables the following inference rule:
\[\text{rdfs:domain}(P,C) \land P(X, Y) \Rightarrow X \in C\]
\emph{Default: ``on''}

	\item[RangeRule] Enables the following inference rule:
\[\text{rdfs:range}(P,C) \land P(X, Y) \Rightarrow Y \in C\]
\emph{Default: ``on''}

	\item[EquivalentClassRule] Enables the following inference rule:
\[\text{owl:equivalentClass}(A,B) \Rightarrow A \subset B \land B \subset A\]
where $\subset$ means ``sub-class of''. \emph{Default: ``on''}

	\item[EquivalentPropRule] Enables the following inference rule:
\[\text{owl:equivalentProperty}(P,Q) \Rightarrow P \sqsubset Q \land Q \sqsubset P\]
where $\sqsubset$ means ``sub-property of''. \emph{Default: ``on''}

	\item[InverseOfRule] Enables the following inference rule:
\[\text{owl:inverseOf}(P, Q) \land P(X, Y) \Rightarrow Q(Y, X)\]
\emph{Default: ``on''}

	\item[SymmetricPropRule] Enables the following inference rule:
\[P \in \text{owl:SymmetricProperty} \land P(X, Y) \Rightarrow P(Y, X)\]
\emph{Default: ``on''}

	\item[FunctionalPropRule] Enables the following inference rule:
\[P \in \text{owl:FP} \land P(Z, X) \land P(Z, Y) \Rightarrow \text{owl:sameAs}(X, Y)\]
where ``FP'' stands for ``FunctionalProperty''. \emph{Default: ``on''}

	\item[InvFunctionalPropRule] Enables the following inference rule:
\[P \in \text{owl:IFP} \land P(X, Z) \land P(Y, Z) \Rightarrow \text{owl:sameAs}(X, Y)\]
where ``IFP'' stands for ``InverseFunctionalProperty''. \emph{Default: ``on''}

	\item[TransitivePropRule] Enables the following inference rule:
\[P \in \text{owl:TransitiveProperty} \land P(X, Y) \land P(Y, Z) \Rightarrow P(X, Z)\]
\emph{Default: ``on''}
\end{description}

%%%%%%%%%%%%%%%%%%%%%%%%%%%%%%%%%%%%%%%%%%%%%%%%%%
\subsection{\acp{pmnt} Logging Configuration}
\label{section-logging-config}

These files configure logging for \acp{pmnt} Java code:
\begin{quote}
	\path|ParliamentKB/config/logback.foreground.xml|\\
	\path|ParliamentKB/config/logback.container.xml|
\end{quote}
They are standard Logback configuration files.  The first applies when running \ac{pmnt} from the command line, and the second is used when running as a container.

The file \path|ParliamentLogConfig.txt| controls native code logging.  Its format is similar to \path|ParliamentKbConfig.txt|, and it is located in a similar way (see Section~\ref{section-main-config}), using file name \path|ParliamentLogConfig.txt| and the environment variable \verb|PARLIAMENT_LOG_CONFIG_PATH|.  The recognized settings are:
\begin{description}
	\item[logToConsole] Indicates whether to log to the console. \emph{Default: ``no''}

	\item[logConsoleAsynchronous] Indicates whether to log asynchronously to the console. \emph{Default: ``no''}

	\item[logConsoleAutoFlush] Indicates whether console log entries should be flushed after each log operation. \emph{Default: ``yes''}

	\item[logToFile] Indicates whether to log to a file. \emph{Default: ``yes''}

	\item[logFilePath] The path of the log file.  May include the following placeholders to add contextual information to the file name:
	\begin{description}[noitemsep]
		\item[\texttt{\%N}]Cardinal number of the log file.  A number between the \% and the N will left-pad with zeros to that many places.
		\item[\texttt{\%Y}]Year
		\item[\texttt{\%m}]Month
		\item[\texttt{\%d}]Day
		\item[\texttt{\%H}]Hours
		\item[\texttt{\%M}]Minutes
		\item[\texttt{\%S}]Seconds
	\end{description}
	\emph{Default:} \texttt{log/ParliamentNative\%3N\_\%Y-\%m-\%d\_\%H-\%M-\%S.log} relative to the kbDirectoryPath setting in \path|ParliamentKbConfig.txt|

	\item[logFileAsynchronous] Indicates whether to log asynchronously to the file. \emph{Default: ``no''}

	\item[logFileAutoFlush] Indicates whether file log entries should be flushed after each log operation. \emph{Default: ``yes''}

	\item[logFileRotationSize] The maximum approximate size (in bytes) of the log file.  When this value is exceeded, the log file will be rotated. \emph{Default: 10 MB}

	\item[logFileMaxAccumSize] The maximum accumulated size of rotated log files.  When this value is exceeded, the oldest log file is deleted. \emph{Default: 150 MB}

	\item[logFileMinFreeSpace] The minimum amount of free space on the volume where log files are stored.  When the free space dips below this value, old log files will be deleted so as to free up enough space. \emph{Default: 100 MB}

	\item[logFileRotationTimePoint] The time of day at which the log file should be rotated regardless of its size.  Must be in the format ``HH:MM:SS''. \emph{Default: 2:00 AM}

	\item[logLevel] The global logging level.  Must be one of TRACE, DEBUG, INFO, WARN, and ERROR. \emph{Default: ``INFO''}

	\item[logChannelLevel] Overrides the global logging level for a single channel.  (Generally a channel is equivalent to a class in the source code.)  The value of this setting takes the form \begin{quote}\emph{«channel-name» = «log-level»}\end{quote} where \emph{«log-level»} takes one of the values listed for the global logging level above.  This option may be repeated to override the level for multiple channels. \emph{Default: None}
\end{description}



%%%%%%%%%%%%%%%%%%%%%%%%%%%%%%%%%%%%%%%%%%%%%%%%%%
\section{\ac{pmnt}-Specific Features}
\label{section-parliament-specific-features}

\ac{pmnt} has a number of features beyond those implied by the \ac{rdf}, \ac{owl}, and \ac{sparql} standards, detailed in the sub-sections below.

%%%%%%%%%%%%%%%%%%%%%%%%%%%%%%%%%%%%%%%%%%%%%%%%%%
\subsection{Reserved Predicates}
\label{section-reserved-predicates}

The \verb|par:directType| and \verb|par:directSubClassOf| predicates are treated specially by \ac{pmnt}.  They are reserved in the sense that they may not be used as the subject, predicate, or object of any  statements inserted in \ac{pmnt}.  However, when used in queries they have special meanings.

In queries, \verb|par:directType| is interpreted like \verb|rdf:type| except that inferred statements are ignored.  Similarly, \verb|par:directSubClassOf| is the same as \verb|rdf:subClassOf| except that it ignores inferred statements.  These predicates provide high-performance shortcuts for queries that are looking for a ``most-derived type.''  For instance, to find the most-derived type of a resource \verb|?x|, we can use the following query:

{\small\begin{verbatim}
prefix rdfs: <http://www.w3.org/2000/01/rdf-schema#>
select distinct ?x ?type where {
   ?x a ?type .
   filter not exists {
      ?x a ?type2 .
      ?type2 rdfs:subClassOf ?type .
      filter ( ?type2 != ?type )
   }
}
\end{verbatim}}

This operates by first finding all of \verb|?x|'s types and then filtering those that have a distinct sub-type.  This is correct and \ac{sparql}-compliant, but the negation is expensive.  Alternately, in \ac{pmnt} we can write this:

{\small\begin{verbatim}
prefix par: <http://parliament.semwebcentral.org/parliament#>
select distinct ?x ?type where {
   ?x par:directType ?type .
}
\end{verbatim}}

This is no longer standards-compliant, and it won't necessarily produce exactly the same results.  However, it is significantly less verbose and, in many cases, much faster.

%%%%%%%%%%%%%%%%%%%%%%%%%%%%%%%%%%%%%%%%%%%%%%%%%%
\subsection{Reification}
\label{section-reification}

\emph{[Yet to be written]}

%%%%%%%%%%%%%%%%%%%%%%%%%%%%%%%%%%%%%%%%%%%%%%%%%%
\subsection{Configuring and Using Indexes}
\label{section-indexes}

\emph{[Yet to be written]}



%%%%%%%%%%%%%%%%%%%%%%%%%%%%%%%%%%%%%%%%%%%%%%%%%%
\section{The ParliamentAdmin Utility}
\label{section-parliamentadmin-utility}

\emph{[Yet to be written]}



%%%%%%%%%%%%%%%%%%%%%%%%%%%%%%%%%%%%%%%%%%%%%%%%%%
\section{Troubleshooting}
\label{section-troubleshooting}

Because \ac{pmnt} is most often used within a Java environment, the most common problem is the dreaded ``Unsatisfied Link Error''.  This error is generated by the \ac{jvm} when the Java code requests the loading of a \ac{dll}, but the \ac{jvm} is unable to load that \ac{dll}.  There are quite a few reasons why you may encounter an unsatisfied link error, and unfortunately, the \ac{jvm} is not very good about generating an illuminating error message.  So, should you encounter this problem, here are some things to keep in mind:
\begin{itemize}
	\item Double-check that you have followed all the deployment instructions correctly for your chosen mode of deployment.

	\item Make sure that your system is up-to-date with respect to patches.  For Windows, this means a visit to the Windows Update web site\urlcite{http://windowsupdate.microsoft.com/}.  Most other operating systems have a similar facility.

	\item Be sure that a 64-bit build of \ac{pmnt} is loaded by a 64-bit version of Java, and that a 32-bit build of \ac{pmnt} is loaded by a 32-bit version of Java.

	\item When the \ac{jvm} loads the \ac{pmnt} \ac{dll}, \ac{pmnt} itself loads several other \acp{dll}.  Even if the \ac{jvm} is able to load \ac{pmnt}, an unsatisfied link error will result if \ac{pmnt} cannot load one of its subsidiary \acp{dll}.  Furthermore, the error message accompanying the unsatisfied link error usually does not indicate which \ac{dll} caused the load failure.  The subsidiary \acp{dll} fall into the following three categories:
	\begin{itemize}
		\item The Boost and RocksDB \acp{dll}:  In the deployment procedures detailed in this document, these \acp{dll} should reside in the same location as \ac{pmnt} itself.
		\item The C and C++ run-time libraries:  On Unix-like systems these should be available automatically.  On Windows you may need to install the Visual Studio run-time libraries.  You can find installers for these in the \path|RedistributablePackages| subdirectory of your binary distribution of \ac{pmnt}.
		\item Various operating system \acp{dll}:  These are rarely a problem, because the OS needs to be able to find these to function.
	\end{itemize}

	\item The Java system property \verb|java.library.path| is a list of paths that Java searches when loading native libraries.  This suggests that all of the operating system-specific rules for locating \acp{dll} can be circumvented by providing a suitable definition for this system property.  Unfortunately, when Java loads a library that itself loads other libraries (as \ac{pmnt} does), Java can consult \verb|java.library.path| only when loading the top-level library, because the loading process for its dependents is managed entirely within the operating system.  Thus, this property does not solve the problem for \ac{pmnt}.

	\item Running \path|ParliamentAdmin| helps to diagnose unsatisfied link errors, because this tool also loads the \ac{pmnt} \ac{dll}, but it does so without involving Java (because \path|ParliamentAdmin| is written in C++).  This technique can be used to isolate the error to either the operating system level (when \path|ParliamentAdmin| fails to load the \ac{dll}) or to the Java level (when \path|ParliamentAdmin| succeeds).  Even the simple command ``\path|ParliamentAdmin -v|'' to get the version requires loading the \ac{pmnt} \ac{dll}, so this is an easy test.

	\item Multiple Java installations can confuse the \ac{dll} loading process.  If you invoke \ac{pmnt} using any of the following:
\begin{quote}
	\path|./parliament|\hfill(on Linux/UNIX)\\
	\path|.\parliament.ps1|\hfill(on Windows)
\end{quote}
	then the Java installation is found by simply invoking the command \path|java| and relying on the operating system to locate the installation.  The \verb|JAVA_HOME| environment variable is not used, unless the operating system uses it.

	\item The process by which \ac{pmnt} loads its \acp{dll} on Windows is a bit complex:

	\begin{itemize}
		\item The \ac{jvm} finds \path|FixupParliamentPath.dll| using the Java system property \verb|java.library.path|.  (This works, because that \ac{dll} has no dependencies of its own.)
		\item Code in \path|FixupParliamentPath.dll| adds the directory where it lives to the path.
		\item Then the \ac{pmnt} \ac{dll} is loaded.  The \ac{jvm} finds it using \verb|java.library.path|, and the operating system finds its dependencies because they, \path|FixupParliamentPath.dll|, and the \ac{pmnt} \ac{dll} all live in the same directory.
	\end{itemize}

	\item On Mac OS X and Linux, the \ac{jvm} finds the \ac{pmnt} shared library by consulting the Java system property \verb|java.library.path|, and the operating system looks for its dependencies via settings placed in the shared libraries during the \ac{pmnt} build process.  The command ``\path|otool -L|'' (on Mac) or \path|ldd| (on Linux) will show from where the operating system is trying to load subsidiary shared libraries.

	\item A debug build of \ac{pmnt} will not work on Windows unless the corresponding version of Visual Studio is installed.  This is because the debug versions of the C and C++ run-time libraries are included only with Visual Studio.

	\item More rarely, unsatisfied link errors can result from a bad build of \ac{pmnt} that causes the symbols exported from the \ac{dll} to be named incorrectly.  If this happen, the \ac{jvm} will succeed in loading the \ac{dll}, but fail to find the entry points it needs.  This results in an unsatisfied link error that is indistinguishable from the cases where the \ac{jvm} cannot load the \ac{dll} at all.  This condition is more likely to occur on Windows, and usually has something to do with the symbol \verb|BUILDING_KBCORE| not being defined on the compiler command line.
\end{itemize}
